\documentclass[11pt]{article}
\usepackage{amsmath,amssymb,amsthm,amsfonts}
\usepackage{graphicx}
\usepackage{algorithm}
\usepackage{algpseudocode}
\usepackage{booktabs}
\usepackage{geometry}
\geometry{margin=1in}

\title{Stochastic Choice of Basis Functions in Adaptive\\Function Approximation and the Functional-Link Net}
\author{Summary of Igelnik and Pao, IEEE Transactions on Neural Networks, 1995}
\date{}

\begin{document}
\maketitle

\section{Introduction}

This paper presents a theoretical justification for the random vector version of the functional-link (RVFL) net based on a general approach to adaptive function approximation. The method uses a limit-integral representation of the function to be approximated and evaluates the integral through the Monte-Carlo approach. The paper provides mathematical proof that RVFL networks are universal approximators for continuous functions and that they achieve approximation error convergence at a rate of $O(C/\sqrt{n})$, where $n$ is the number of basis functions, with $C$ independent of $n$.

\section{Methodology}

\subsection{General Approach to Function Approximation}

The key insight of the approach is to represent a continuous function $f \in C(I^d)$ defined on the standard hypercube $I^d = [0,1]^d \subset \mathbb{R}^d$ as a limit-integral:

\begin{align}
f(x) = \lim_{\lambda \rightarrow a} \int_V G_{\lambda,x}(X)T[f(X)]dX
\end{align}

where:
\begin{itemize}
    \item $\lambda$ and $X$ are parameters (with $\lambda$ being low-dimensional and $X$ multidimensional)
    \item $G$ is an activation function 
    \item $T$ is an operator defined on $C(I^d)$
    \item $a$ is a finite or infinite real number
    \item $V$ is the domain of parameter $X$
\end{itemize}

The approximation is achieved through two stages:
\begin{enumerate}
    \item Approximating the limiting value of the integral by the integral with finite but large $\lambda$
    \item Obtaining an estimate of the multiple integral using the Monte-Carlo method:
    \begin{align}
    \int_V G_{\lambda,x}(X)T[f(X)]dX \approx \frac{|V|}{n}\sum_{k=1}^n G_{\lambda,x}(X_k)T[f(X_k)]
    \end{align}
\end{enumerate}

where $X_1, \ldots, X_n$ is a sample of size $n$ drawn from the uniform distribution on $V$.

\subsection{Random Vector Functional-Link (RVFL) Net}

The RVFL net is defined as a single hidden layer feedforward neural network:

\begin{align}
f_{w_n}(x) = \sum_{k=1}^n a_k g(w_k \cdot x + b_k)
\end{align}

where:
\begin{itemize}
    \item $w_n = (n, a_1, \ldots, a_n, b_1, \ldots, b_n, w_1, \ldots, w_n)$ is the overall parameter of the net
    \item $w_k \cdot x$ is the inner product of vectors $w_k$ and $x$
    \item $g$ is an activation function
    \item $a_k$ are trainable output layer weights
    \item $w_k \in \mathbb{R}^d$ and $b_k \in \mathbb{R}$ are randomly assigned and fixed hidden layer parameters
\end{itemize}

The key distinction between RVFL and general nonlinear perceptron (GNP) is:
\begin{itemize}
    \item In RVFL: Parameters $w_k, b_k$ of the hidden layer are selected randomly and independently in advance, while parameters $a_k$ of the output layer are learned using simple quadratic optimization
    \item In GNP: All parameters $a_k, w_k, b_k$ need to be learned using complex nonquadratic optimization
\end{itemize}

\subsection{Algorithm Details}

The RVFL learning algorithm proceeds as follows:

\begin{algorithm}
\caption{RVFL Learning Algorithm}
\begin{algorithmic}[1]
\State Generate random parameters $w_k \in \mathbb{R}^d$ and $b_k \in \mathbb{R}$ for $k = 1,\ldots,n$
\State For $w_k$: Generate independently and uniformly from $[0,\Omega] \times [-\Omega,\Omega]^{d-1}$
\State For $b_k$: Set $b_k = -w_k \cdot y_k - u_k$ where $y_k$ is uniformly distributed in $I^d$ and $u_k$ is uniformly distributed in $[-2\Omega,2\Omega]$
\State Compute the output of hidden nodes for all training samples
\State Find the optimal values for output weights $a_k$ using linear least squares optimization
\end{algorithmic}
\end{algorithm}

The network training is thus reduced to linear optimization (typically using conjugate gradient method), making it simple and fast.

\section{Theoretical Findings}

The paper presents three main theoretical results:

\subsection{Universal Approximation Property}

\textbf{Theorem 1:} For any compact set $K \subset I^d$, $K \neq I^d$ and any absolutely integrable activation function $g$ such that $\int_{\mathbb{R}} g^2(x) dx < \infty$, there exists a sequence of RVFL networks $\{f_{w_n}\}$ and a sequence of probability measures $\{p_{n,\Omega,\alpha}\}$ such that:

\begin{align}
E\int_K |f(x) - f_{w_n}(x)|^2 dx \rightarrow 0 \text{ as } n \rightarrow \infty
\end{align}

\textbf{Corollary 1:} The universal approximation property also holds for any differentiable activation function $g$ with $\int_{\mathbb{R}} [g'(x)]^2 dx < \infty$, which includes commonly used sigmoid functions.

\textbf{Corollary 2:} An alternative form of the RVFL network as a product of univariate functions:

\begin{align}
f_{w_n}(x) = \sum_{k=1}^n a_k \prod_{i=1}^d g(w_{ki}x_i+b_{ki})
\end{align}

also possesses the universal approximation property.

\subsection{Radial Basis Function Networks with Random Parameters}

\textbf{Theorem 2:} Neural networks with radial basis functions and random parameters in the form:

\begin{align}
f_{w_n}(x) = \sum_{k=1}^n a_k g\left(\sum_{i=1}^d w_{ki}(x_i-y_{ki})^2\right)
\end{align}

are also universal approximators for continuous functions on any compact subset of $I^d$.

\subsection{Approximation Error Rate}

\textbf{Theorem 3:} For any Lipschitz continuous function $f \in C(I^d)$ satisfying $|f(x) - f(y)| \leq K\|x-y\|$ for some constant $K > 0$, any compact set $K \subset I^d$, $K \neq I^d$, and any activation function $g_{\beta}$ with appropriate support, there exists a sequence of RVFL networks $\{f_{w_n}\}$, a sequence of probability measures $\{p_{n,\Omega,\alpha}\}$, and a constant $C_{f,g,\Omega,\alpha,\beta,d}$ such that:

\begin{align}
E\int_K |f(x) - f_{w_n}(x)|^2 dx \leq \frac{C_{f,g,\Omega,\alpha,\beta,d}}{n}
\end{align}

This means the approximation error converges to zero at a rate of $O(C/\sqrt{n})$ with $C$ independent of $n$. This is significantly better than the approximation by linear combinations of fixed basis functions, which gives an error on the order of $O(1/n^{1/d})$.

\section{Proof Sketch}

The proofs rely on several key steps:

\begin{enumerate}
    \item Representing a function to be approximated as the limiting value of a multidimensional integral using window transformations
    \item Using a kernel that concentrates a function at specific points or surfaces of its domain
    \item Replacing the kernel by the activation function $g$
    \item Using Monte-Carlo methods to approximate the integral
    \item Establishing bounds on the approximation error
\end{enumerate}

The mathematical foundation relies on representing the function using a window transformation $h_x(y,w)$ that approaches the multidimensional delta function as parameters approach infinity. This leads to a limit-integral representation which is then evaluated using Monte-Carlo integration.

\section{Numerical Examples and Applications}

The paper describes several practical applications where RVFL networks demonstrated superior performance compared to traditional neural networks:

\subsection{Optical Ellipsometry Problem}

In molecular beam epitaxy film thickness monitoring, RVFL networks successfully solved the complex inversion problem of determining film refractive index and thickness from ellipsometry measurements:

\begin{itemize}
    \item Input variables: Two pairs of $\psi$ and $\Delta$ angular measurements
    \item Output variables: $n$ and $\kappa$ (real and imaginary parts of refractive index) and two film thicknesses $d_0$ and $d_1$
    \item Training: $10^4$ training patterns (sparse set in 4D space)
    \item Results: Refractive index $n$ evaluated with error less than 0.1\%
    \item Performance: Training achieved in about six hours on a SPARC 2 workstation, whereas backpropagation failed to achieve satisfactory results
    \item Network size: 50-200 basis functions
\end{itemize}

\subsection{Comparison with Backpropagation}

A comparison table shows RVFL outperforming backpropagation (BP) across various tasks:

\begin{center}
\begin{tabular}{lccccc}
\toprule
\textbf{Task} & \textbf{Dim.} & \textbf{Size} & \textbf{RVFL Time} & \textbf{BP Time} & \textbf{Accuracy} \\
\midrule
Pattern Rec. & 25-5 & 100 & 46 sec & 31.2 min & Equal \\
Image Gen. & 2-32 & 400 & 165 sec & 240.8 min & Equal \\
Ellipsometry & 4-4 & 9000 & 6 hrs & * & RVFL better \\
\bottomrule
\end{tabular}
\end{center}
* Backpropagation failed to achieve comparable accuracy

\section{Implications and Future Directions}

The theoretical and practical results presented have several important implications:

\subsection{Computational Efficiency}

RVFL networks offer significant advantages over traditional neural networks:

\begin{itemize}
    \item Linear learning (only output weights are trained) is simple and fast
    \item Approximation error converges at a rate of $O(C/\sqrt{n})$, which is the same as GNP but with much simpler training
    \item RVFL avoids the exponential growth in the number of basis functions that occurs with fixed basis functions
\end{itemize}

\subsection{Enhancements for Improved Accuracy}

The paper suggests methods to enhance the accuracy of RVFL networks:

\begin{itemize}
    \item Using ensemble methods with multiple RVFL networks
    \item Applying the Pincus formula for estimating the optimal parameters:
    \begin{align}
    \theta_j \approx \frac{\sum_{k=1}^K \theta_{kj} \exp[-\lambda E_{n,N}(\theta_k)]}{\sum_{k=1}^K \exp[-\lambda E_{n,N}(\theta_k)]}, j=1,\ldots,2n
    \end{align}
    where $\theta_k = (\theta_{k1},\ldots,\theta_{k,2n})$ are randomly generated parameter sets
    \item Using variance reduction methods in Monte-Carlo integration where specific information about the function is available
\end{itemize}

\subsection{Broader Impact}

The stochastic approach based on limit-integral representation with Monte-Carlo evaluation opens new possibilities for efficient approximation of multivariate functions:

\begin{itemize}
    \item The approach can be extended to other types of neural networks
    \item It provides a mathematical foundation for understanding why randomized networks can be effective
    \item It suggests that combining multiple independent RVFL estimates can be more effective than increasing the number of basis functions in a single network
\end{itemize}

The theoretical results confirm that RVFL networks represent a practical compromise between fully trainable networks and networks with fixed basis functions, combining both simplicity of learning and efficiency of representation.

\end{document} 